% ===========================================================================
% Methods section, \input in main.tex under \section{Methods}.
% Every number from the main-study data is a macro from out/methods_numbers.tex,
% written by analysis/R/04_methods_numbers.R: rerun that script, do not type numbers here.
% Pilot numbers (n, item counts) come from the pilot data in the Qualtrics project.
% Item wording: appendix_items.tex.
% ===========================================================================

\subsection{Study design}

The study is an observational, cross-sectional online survey of Dutch adults. It was not
preregistered. The instrument was developed over two pilot studies and then fielded once in a sample
drawn from the Kieskompas panel. The study was approved by the Ethics Review Board of the Psychology
research institute, University of Amsterdam (FMG-20487, 22 August 2026).

The study introduces a new measure of drug policy preference, a nine-point regulation scale asked per
substance, and a first set of attitude items intended to explain positions on it.

% ---------------------------------------------------------------------------
\subsection{Development of the instrument}

\subsubsection{The regulation scale}

Drug policy is multidimensional, varying among other things in who may supply a substance, to whom,
on what terms, and how strictly rules are enforced. Following \citet{maccoun2001}, we take the
restrictiveness of legal access as the principal dimension and order the answer options along it.
This yields nine regulatory models, each defined in one or two sentences of plain Dutch
(Figure~\ref{fig:policy_scale}): a ban (supply, purchase and possession prohibited); research only;
treatment only (available from a doctor to people who are seriously ill or addicted); on prescription
(a doctor may prescribe it to anyone who would benefit); sale in a specialist shop, without and with
profit; sale by ordinary shops and hospitality, without and with advertising; and unrestricted sale.

The scale deliberately leaves out a separate option for depenalised use. In \citeauthor{maccoun2001}'s
analysis depenalisation differs from prohibition in the severity of sanctions, not in the
restrictiveness of access, and so lies on a second dimension. Sanctions for personal possession are
measured instead by a belief item (Table~\ref{tab:measures}).

Respondents rated nine substances: alcohol, tobacco and vapes, cannabis, psychedelics (such as
mushrooms and LSD), MDMA, cocaine/crack, speed (amphetamine), heroin and strong painkillers. For each
substance the nine options were shown as nine ordered buttons, with a separate \emph{no opinion}
option. Substances were presented in random order.

Before rating the substances, respondents answered two practice items about the current Dutch rules
for two substances outside the study: ADHD medication (on prescription only) and coffee (sold
everywhere, advertising allowed). A correct answer was required to continue, so that every respondent
had used the scale once before the substances.

\subsubsection{The attitude items}

\textbf{Theoretical basis.} Items were written to cover the three classical components of an
attitude: cognition (beliefs), affect (feelings) and behaviour \parencite{rosenberg1960,eagly1993}.
How the items were selected follows from the Causal Attitude Network (CAN) model
\parencite{dalege2016,dalege2018}. CAN treats an attitude as a network of evaluative reactions that
influence one another. The aim of item construction was therefore to map the reactions that plausibly
make up a person's stance on drug regulation, and to measure each with one item. Following
recommendations for network studies \parencite{burger2022,chambon2026}, we wrote items with distinct
content.

\textbf{Content.} Belief items cover four domains: \emph{addiction}, \emph{crime}, the consequences of
\emph{regulation}, and \emph{society}, the last covering norms about autonomy, profit from addictive
substances and the international prohibition \parencite[cf.][]{vangelder2024}. Emotion items record
whether respondents feel anger, fear or worry, or joy about six topics. Behaviour items describe what
respondents would do in situations involving drugs. Own use and involvement with the topic were
measured separately.

\textbf{Wording.} Item wording drew on \citeauthor{saris2004}'s \citeyearpar{saris2004} typology of
concepts-by-intuition, which links the concept an item measures to the linguistic form of the
assertion. Where possible we preferred consequence beliefs (``If drugs could be sold legally, \ldots'')
and statements about specific interventions to global evaluations (``The war on drugs is a failure'').
Global evaluations predict policy stance well \parencite{bergquist2022}, but partly because they are
conceptually close to the outcome \parencite{vysevkova2026}. All items were written by the authors in
Dutch, at B1 level. The English wordings in Appendix~\ref{app:items} are our own translations.

\subsubsection{Pilot studies}

Two pilots were run on earlier versions of the instrument: a convenience sample ($n = 121$, August
2026) and a quota sample of Dutch adults recruited through Prolific ($n = 120$, September 2026). In
both, policy preference was measured on a 0--100 slider with seven anchors. Pilot 1 had 55 attitude
items and pilot 2 had 47.

Items were screened for redundancy (Spearman $\rho$ with other items), for their contribution to
explaining the policy sliders (leave-one-out $\Delta R^2_{\mathrm{adj}}$), and for their corrected
item--rest correlation within the block \parencite{burger2022,chambon2026}. On the basis of these
results, items that added little were dropped or reworded, and own use was asked per substance group
rather than per substance. The slider was replaced by the nine ordered options described above, which
name each regulatory model explicitly rather than placing it between anchors. The main study has 48
attitude items (Table~\ref{tab:measures}).

% ---------------------------------------------------------------------------
\subsection{Main study}

\subsubsection{Participants and procedure}

The main study was fielded through the VIP panel of Kieskompas from \fieldStart\ to \fieldEnd. The VIP
panel is a stratified random sample from Kieskompas's opt-in Grote Burgerpanel, stratified on Nielsen
region, sex, age, education and migration background, with a correction for vote recall; the
population frame is the CBS Gouden Standaard. Respondents had to be 18 or older and give informed
consent.

Of the \nInvited\ panel members invited, \nStarted\ started the survey and \nCompleted\ completed it
(response rate \responseRate\%). The \nBrokeOff\ who broke off were excluded. Of the completed
interviews, \nAnalytic\ provided the background variables needed for weighting and form the analytic
sample. Median completion time was \medianMinutes\ minutes.

Kieskompas weighted the sample by iterative post-stratification to marginal distributions from the CBS
Gouden Standaard of age (four categories), sex (two), education (three), migration background
(three: Dutch, EU, non-EU) and Nielsen region (five), and to the result of the 29 October 2025
parliamentary election, with each party in parliament weighted separately and the remaining parties,
blank votes and non-voters as further categories. Weights were trimmed at the 99th percentile. All
estimates are weighted unless stated otherwise. Table~\ref{tab:sample_composition} gives the unweighted and weighted
sample composition.

\begin{table}[htbp]
\centering\small
\caption{Sample composition ($n = 1{,}239$): number of respondents, unweighted and weighted percentages. Region in Nielsen regions; urbanisation in CBS address-density classes.}
\label{tab:sample_composition}
\begin{tabular}{@{}lrrr@{}}
\toprule
 & $n$ & Unweighted \% & Weighted \% \\
\midrule
\addlinespace
\multicolumn{4}{@{}l}{\textit{Age}} \\
\quad 18--34 & 259 & 20.9 & 28.2 \\
\quad 35--49 & 266 & 21.5 & 21.4 \\
\quad 50--64 & 390 & 31.5 & 25.5 \\
\quad 65+ & 324 & 26.2 & 24.9 \\
\addlinespace
\multicolumn{4}{@{}l}{\textit{Sex}} \\
\quad Male & 674 & 54.4 & 49.4 \\
\quad Female & 565 & 45.6 & 50.6 \\
\addlinespace
\multicolumn{4}{@{}l}{\textit{Education}} \\
\quad Theoretical & 512 & 41.3 & 32.4 \\
\quad Middle & 488 & 39.4 & 43.2 \\
\quad Practical & 239 & 19.3 & 24.4 \\
\addlinespace
\multicolumn{4}{@{}l}{\textit{Region}} \\
\quad Three largest cities and suburbs & 164 & 13.2 & 16.5 \\
\quad Rest of the west & 382 & 30.8 & 30.0 \\
\quad North & 139 & 11.2 & 10.4 \\
\quad East & 263 & 21.2 & 21.6 \\
\quad South & 291 & 23.5 & 21.6 \\
\addlinespace
\multicolumn{4}{@{}l}{\textit{Urbanisation}} \\
\quad Very high & 315 & 25.4 & 28.4 \\
\quad High & 392 & 31.6 & 31.1 \\
\quad Moderate & 198 & 16.0 & 15.0 \\
\quad Low & 234 & 18.9 & 17.3 \\
\quad Not urban & 77 & 6.2 & 5.6 \\
\quad Unknown & 23 & 1.9 & 2.7 \\
\addlinespace
\multicolumn{4}{@{}l}{\textit{Vote recall, 2025 election}} \\
\quad GroenLinks-PvdA & 240 & 19.4 & 10.8 \\
\quad D66 & 234 & 18.9 & 14.0 \\
\quad VVD & 133 & 10.7 & 11.6 \\
\quad PVV & 125 & 10.1 & 13.4 \\
\quad CDA & 114 & 9.2 & 9.6 \\
\quad JA21 & 79 & 6.4 & 4.9 \\
\quad SP & 42 & 3.4 & 1.6 \\
\quad FvD & 37 & 3.0 & 3.7 \\
\quad PvdD & 36 & 2.9 & 1.7 \\
\quad BBB & 33 & 2.7 & 2.2 \\
\quad Volt & 28 & 2.3 & 1.0 \\
\quad SGP & 20 & 1.6 & 1.8 \\
\quad 50PLUS & 19 & 1.5 & 1.2 \\
\quad CU & 15 & 1.2 & 1.5 \\
\quad DENK & 1 & 0.1 & 0.5 \\
\quad Other & 19 & 1.5 & 1.2 \\
\quad Blank & 5 & 0.4 & 0.3 \\
\quad Did not vote & 59 & 4.8 & 19.1 \\
\bottomrule
\end{tabular}
\end{table}


The target sample size was 1{,}000. The weighting gives a design effect of \designEffect, so the
effective sample size is \nEff\ \parencite{kish1965}, and a share of 50\% is estimated with a 95\%
margin of error of $\pm$\marginError\ percentage points.

\subsubsection{Measures}

The questionnaire was administered in Dutch, in this order: policy preference, beliefs, behaviour,
own use, emotions, involvement (Table~\ref{tab:measures}). Background variables (age, sex, education,
region, urbanisation and vote recall) came from the panel profile. Belief and behaviour items were
answered on seven-point agreement scales (\emph{strongly disagree} to \emph{strongly agree}); belief
items had a \emph{no opinion} option. Items on which agreement favours prohibition were
reverse-scored, so that high scores indicate a more liberal position. \emph{No opinion} was chosen
for \noOpPolicyMin--\noOpPolicyMax\% of the policy ratings, depending on the substance, and for at
most \noOpBeliefMax\% of answers to any one belief item; it is treated as missing. The own-use items
stated that answers were anonymous and could be skipped. The survey closed with a referral to the
Drugs Infolijn, the national drug information line.

\begin{table}[htbp]
\centering
\small
\caption{Measures in the main study. Full item wording in Appendix~\ref{app:items}.}
\label{tab:measures}
\begin{tabularx}{\linewidth}{@{}l l c l X@{}}
\toprule
Block & Domain & Items & Format & Example \\
\midrule
Policy preference & -- & 9 & 9 options + no opinion & one rating per substance \\
\addlinespace
Beliefs (cognition) & Society & 5 & 7-point + no opinion & People should be able to decide for themselves whether to use harmful substances. \\
 & Addiction & 5 & 7-point + no opinion & If a substance can be bought legally, more people become addicted to it. \\
 & Regulation & 4 & 7-point + no opinion & Allowing the sale of drugs would earn the Netherlands a lot of money. \\
 & Crime & 4 & 7-point + no opinion & Possessing drugs for personal use should not be a criminal offence. \\
\addlinespace
Behaviour & -- & 3 & 7-point & I avoid people who use drugs. \\
\addlinespace
Emotions (affect) & 6 topics & 18 & check all that apply & anger / fear or worry / joy, for alcohol, tobacco, soft drugs, hard drugs, the war on drugs, legalisation \\
\addlinespace
Own use & -- & 4 & 4 categories & never \ldots\ often in the past year; alcohol, tobacco/vapes, soft drugs, hard drugs \\
\addlinespace
Involvement & -- & 5 & slider 0--100 & How a party deals with drugs and drug crime affects my vote. \\
\bottomrule
\end{tabularx}
\end{table}

% ---------------------------------------------------------------------------
\subsection{Statistical analysis}

All analyses were done in R \rVersion\ \parencite{rcore}, using the \texttt{survey} package
\parencite{lumley2004} for weighted estimates and design-based standard errors. Code is available at
\url{https://github.com/timovanommeren/Measuring_drug_policy_as_a_spectrum}; the data are available
from the authors on request.

\textbf{Distribution of preferences.} For each substance we report the weighted share choosing each of
the nine options, with design-based 95\% confidence intervals. We compare shares across age,
education, urbanisation and political orientation, and summarise each party's voters by their mean
position on the illicit and on the legal substances. That summary treats the nine options as equally
spaced.

\textbf{Explaining preferences.} Each substance's policy preference is regressed on the attitude
items, using the survey weights and design-based standard errors. [TODO: linear or ordinal model;
covariates.] We report the variance explained per block and per substance.
